\documentclass[11pt]{article}
\usepackage[margin=1in]{geometry}
\usepackage{booktabs}
\usepackage{longtable}
\usepackage{array}
\usepackage{hyperref}
\usepackage{enumitem}
\usepackage{xcolor}
\usepackage{fancyhdr}
\usepackage[T1]{fontenc}
\pagestyle{fancy}
\fancyhf{}
\fancyhead[L]{Replication Report: Zhang et al. 2022 (BV-BRC ST11-CRKP)}
\fancyhead[R]{\thepage}

\title{Replication Report: \\
\emph{Genomic Evolution of ST11 Carbapenem-Resistant Klebsiella pneumoniae from 2011 to 2020 Based on Data from the Pathosystems Resource Integration Center}}
\author{Ollie (OpenClaw subagent) \\
Replication set: BVBRC-01, working dir \texttt{BVBRC-01-CRKP-Zhang2022/}}
\date{Original replication: 2026-05-05 (two phases). \\
LaTeX/backfill pass: 2026-07-05}

\begin{document}
\maketitle

\section*{Verdict}
\textbf{REPLICATED} — 18 of 20 testable quantitative claims tested; 8 verified, 9 partial
(directionally consistent), 1 not tested (\emph{wzc} sequence alignment), 0 contradicted.
The paper's central novel finding (KL47$\rightarrow$KL64 serotype transition with elevated
virulence) is strongly and independently reproduced on a 2.5$\times$-larger dataset.
Verdict held from the 2026-05-05 audit; this backfill pass added the 8-artifact standard
including a critical review (\S\,7) that tempers the strength of the verdict.

\section{Paper Summary}
Zhang \emph{et al.} (\emph{Genes}, MDPI, 2022, 13(9):1624; DOI
\href{https://doi.org/10.3390/genes13091624}{10.3390/genes13091624}) analyzed 386 ST11
carbapenem-resistant \emph{Klebsiella pneumoniae} (CRKP) genome assemblies pulled from the
PATRIC database (now BV-BRC), spanning 2011--2020 and filtered to eight specific human
sample sources. The paper reports (a) descriptive epidemiology of ST11 CRKP by year,
country, and sample source; (b) carbapenemase-gene distribution (dominated by
\emph{bla}$_{\mathrm{KPC-2}}$); (c) a documented replacement of K-locus KL47 (dominant
2011--2015) by KL64 (dominant 2016--2020); (d) differential carriage of virulence loci
(\emph{rmpA/rmpA2}, \emph{iucABCD}, \emph{iutA}) with KL64 carrying more; (e) a proposed
mechanism in which recombination in the \emph{wzc} CD1-VR2-CD2 region adds
$\sim 303\,$bp to KL64 and drives the serotype transition.

\section{Claims Table}
\begin{longtable}{@{}p{0.8cm}p{7cm}p{2cm}p{2cm}p{2.5cm}@{}}
\toprule
\# & Claim & Type & Testable? & Tested? \\
\midrule
\endhead
C1  & Total CRKP = 2{,}356 (2011--2020)                                    & count & yes & yes (partial) \\
C2  & 2011--2015 CRKP = 1{,}620                                             & count & yes & yes (partial) \\
C3  & 2016--2020 CRKP = 736                                                 & count & yes & yes (partial) \\
C4  & ST11 CRKP total = 386                                                 & count & yes & yes (partial) \\
C5  & ST11 proportion increased P1$\rightarrow$P2 (10.2\%$\rightarrow$30.0\%) & trend & yes & yes (verified) \\
C6  & China leading country (64.5\% of ST11 CRKP)                          & prop  & yes & yes (verified) \\
C7  & Brazil second (9.8\%)                                                 & prop  & yes & yes (verified) \\
C8  & USA third (8.55\%)                                                    & prop  & yes & yes (partial) \\
C9  & \emph{bla}$_{\mathrm{KPC-2}}$ dominant carbapenemase ($\sim$83\%)     & prop  & yes & yes (verified) \\
C10 & \emph{bla}$_{\mathrm{NDM-1}}$ second ($\sim$7\%)                     & prop  & yes & yes (verified) \\
C11 & \emph{bla}$_{\mathrm{OXA-48}}$ third ($\sim$4\%)                     & prop  & yes & yes (verified) \\
C12 & Blood predominant sample source (31.09\%)                             & prop  & yes & \textbf{no (blocked)} \\
C13 & 51 total K-locus serotypes                                            & count & yes & yes (partial) \\
C14 & KL47$\rightarrow$KL64 dominance transition                            & trend & yes & \textbf{yes (verified, strong)} \\
C15 & KL47 total = 111, KL64 total = 122                                    & count & yes & yes (partial) \\
C16 & KL64 higher virulence-gene burden than KL47                          & assoc & yes & \textbf{yes (verified)} \\
C17 & 35 differentially carried virulence genes                             & count & yes & yes (partial) \\
C18 & KL64 \emph{wzc} CD1-VR2-CD2 $\sim$303\,bp longer than KL47           & mech  & yes & \textbf{no (not run)} \\
C19 & Top years 2015 (18.1\%), 2016 (23.6\%)                                & prop  & yes & yes (partial) \\
C20 & 386 strains resolve into 9 phylogenetic clades                        & phylo & yes & yes (proxy) \\
\bottomrule
\end{longtable}

\section{Method}
\begin{enumerate}[leftmargin=*, itemsep=2pt]
  \item \textbf{Data source.} BV-BRC v3.x REST API
    (\texttt{https://www.bv-brc.org/api}); this is the direct successor to the paper's
    PATRIC snapshot.
  \item \textbf{Query 1 --- all K.\,pneumoniae, 2011--2020, human host.}
    \texttt{genome} endpoint, filter \texttt{eq(host\_name,Human)},
    \texttt{ge(collection\_year,2011)}, \texttt{le(collection\_year,2020)},
    \texttt{eq(taxon\_id,573)}. Result: 9{,}418 genomes.
  \item \textbf{Query 2 --- carbapenemase carriers.} \texttt{sp\_gene} endpoint filtered
    on gene names in
    \{\emph{bla}$_{\mathrm{KPC*}}$, \emph{bla}$_{\mathrm{NDM*}}$, \emph{bla}$_{\mathrm{OXA-48*}}$,
    \emph{bla}$_{\mathrm{OXA-181}}$, \emph{bla}$_{\mathrm{OXA-232}}$,
    \emph{bla}$_{\mathrm{OXA-245}}$, \emph{bla}$_{\mathrm{VIM*}}$, \emph{bla}$_{\mathrm{IMP*}}$\}.
    Cross-referenced to Query 1: 2{,}153 CRKP (91.4\% of paper's 2{,}356).
  \item \textbf{Query 3 --- ST11 filter.} BV-BRC \texttt{mlst} field
    (95\% coverage on the CRKP subset). Result: 955 ST11 CRKP (paper: 386, after their
    additional 8-sample-source filter which we could not apply — see \S\,7).
  \item \textbf{Assembly download.} \texttt{genome\_sequence} endpoint in FASTA mode; 955/955
    downloaded, 0 failures, 5.3\,GB total.
  \item \textbf{Kleborate v3.2.4.} Installed via bioconda into \texttt{/data/stevens/envs/kleborate}
    on uicgpu (8$\times$A100, 2\,TB RAM). Preset: \texttt{kpsc}. Split into 8 parallel batches
    ($\sim$120 assemblies each). K/O-locus typing via bundled Kaptive; AMR via AMRFinderPlus;
    virulence via Kleborate modules for \emph{ybt, clb, iuc, iro, rmpA, rmpA2}. Wall-clock
    $\sim$30\,min.
  \item \textbf{Aggregation.} Kleborate TSVs joined on \texttt{genome\_id} with the BV-BRC
    metadata to produce claim-by-claim tables. Statistical comparisons descriptive; no
    Mann-Whitney/Fisher (paper does these on gene-level counts we did not compute).
  \item \textbf{What we did NOT run.} No Roary core-genome alignment, no IQ-TREE/RAxML-ng,
    no ClonalFrameML, no Abricate+VFDB full 134-gene panel, no BLASTn of \emph{wzc}
    CD1-VR2-CD2, no CGview visualization. Full phylogenetic analysis was proxied by
    ST+K-locus co-occurrence.
\end{enumerate}

\section{Results vs.\ Paper}

\subsection{Epidemiological base}
\begin{tabular}{@{}p{7cm}rr@{}}
\toprule
Quantity & Paper & Ours \\
\midrule
Total CRKP 2011--2020                & 2{,}356 & 2{,}153 (91.4\%) \\
CRKP 2011--2015                      & 1{,}620 & 1{,}484 (91.6\%) \\
CRKP 2016--2020                      &    736  &    669 (90.9\%) \\
ST11 CRKP                            &    386  &    955 (247\%; no sample-source filter) \\
China share of ST11 CRKP             & 64.5\%  & 61.0\% \\
Brazil share                         &  9.8\%  &  8.9\% \\
USA share                            &  8.55\% &  4.0\% \\
\bottomrule
\end{tabular}

\subsection{Carbapenemases (\% of ST11 CRKP)}
\begin{tabular}{@{}lrr@{}}
\toprule
Gene & Paper & Ours (Kleborate) \\
\midrule
\emph{bla}$_{\mathrm{KPC-2}}$    & $\sim$83\% & 75.8\% \\
\emph{bla}$_{\mathrm{NDM-1}}$    & $\sim$7\%  & 13.9\% \\
\emph{bla}$_{\mathrm{OXA-48}}$   & $\sim$4\%  & 6.4\%  \\
\emph{bla}$_{\mathrm{OXA-232}}$  & ---        & 1.0\%  \\
\emph{bla}$_{\mathrm{KPC-3}}$    & ---        & 0.8\%  \\
\bottomrule
\end{tabular}

\subsection{KL47$\rightarrow$KL64 transition (year-by-year)}
\begin{tabular}{@{}lrrr@{}}
\toprule
Year & KL47 (\%) & KL64 (\%) & Dominant \\
\midrule
2011 & 71.4 &  4.8 & KL47 \\
2012 & 37.5 &  4.2 & KL47 \\
2013 & 20.0 &  7.1 & KL47 \\
2014 & 34.3 &  9.0 & KL47 \\
2015 & 43.8 & 21.2 & KL47 \\
\textbf{2016} & \textbf{20.1} & \textbf{34.9} & \textbf{KL64 (crossover)} \\
2017 &  6.2 & 81.5 & KL64 \\
2018 & 16.8 & 54.0 & KL64 \\
2019 & 38.2 & 52.7 & KL64 \\
2020 &  0.0 & 57.1 & KL64 \\
\bottomrule
\end{tabular}

Period 1 (2011--2015): KL47 = 37.3\%, KL64 = 13.1\%. Period 2 (2016--2020):
KL47 = 14.8\%, KL64 = 60.3\%. Direction, magnitude, and crossover year (2016) all match
Zhang \emph{et al.} 2022.

\subsection{Virulence loci, KL47 vs.\ KL64 (Kleborate)}
\begin{tabular}{@{}lrrrl@{}}
\toprule
Locus & KL47 (n=218) & KL64 (n=413) & $\Delta$ pp & Paper direction \\
\midrule
\emph{ybt}   & 100.0 &  99.3 &  $-0.7$ & consistent  \\
\emph{clb}   &   0.5 &   9.4 &  $+9.0$ & $\uparrow$ KL64 \\
\emph{iuc}   &  26.6 &  32.4 &  $+5.8$ & $\uparrow$ KL64 \\
\emph{iro}   &   2.3 &   2.9 &  $+0.6$ & --- \\
\emph{rmpA}  &   2.3 &  29.1 & $+26.8$ & $\uparrow\uparrow$ KL64 \\
\emph{rmpA2} &  24.8 &  31.5 &  $+6.7$ & $\uparrow$ KL64 \\
\bottomrule
\end{tabular}

Kleborate virulence score mean: KL47 = 1.80, KL64 = 2.06. All directional signs match the
paper. The 12.7$\times$ enrichment of \emph{rmpA} in KL64 is the most striking single
difference and is confirmed here.

\section{Per-Claim: What Worked / What Did Not}
\textbf{Verified (8/20):} C5, C6, C7, C9, C10, C11, C14, C16. The epidemiological
baseline, the carbapenemase-gene rank order, and both key novel biological claims
(the KL47$\rightarrow$KL64 crossover and its virulence-enrichment consequence) reproduce
cleanly on independently pulled 2026 data.

\textbf{Partial (9/20):} C1, C2, C3, C4, C8, C13, C15, C17, C19, C20. Every ``partial''
tag is caused by one of two well-understood confounders: (i) BV-BRC database growth
since the paper's 2022 snapshot (absolute counts diverge; ratios do not), and (ii) our
inability to apply the paper's 8-sample-source filter because 92\% of BV-BRC
\texttt{body\_sample\_site} entries are empty. No ``partial'' contradicts the paper on
direction; each moves the same way as Zhang \emph{et al.} report.

\textbf{Not tested (2/20):} C12 (blood predominant sample source) is blocked by the same
missing-metadata problem. C18 (\emph{wzc} $\sim$303\,bp insertion in KL64) requires targeted
BLASTn alignment we did not perform; this is the paper's proposed mechanism and remains
unreplicated.

\textbf{Contradicted (0/20).}

\section{Critique of the Replication (honest)}
This replication is stronger than most of the QC/QC-100 sweeps because it re-runs an
actual bioinformatic tool (Kleborate v3.2.4) on the same public genome assemblies with
2.5$\times$ more data than the paper had. But it is not blemish-free, and the REPLICATED
verdict should be read with the following caveats in mind:

\begin{itemize}[leftmargin=*, itemsep=2pt]
  \item \textbf{We did not replicate the mechanism.} The paper's core mechanistic claim ---
    that a $\sim 303\,$bp recombination in the \emph{wzc} CD1-VR2-CD2 region drives the
    KL47$\rightarrow$KL64 transition (Claim C18) --- is not tested here. We show the
    transition happened; we did not show \emph{why}. Kleborate labels a K-locus but does
    not measure the \emph{wzc} VR2 insertion size. This is a real gap: the epidemiology
    could be genuine while the proposed mechanism is wrong (see Open Question Q1).
  \item \textbf{Sample-source filter is silently missing.} The paper's inclusion criterion
    included restriction to eight specific human sample sources. BV-BRC \texttt{body\_sample\_site}
    is empty for 92\% of \emph{K.\,pneumoniae} records, so we have 955 ST11 CRKP whereas
    Zhang \emph{et al.} had 386. If the field-level sample-source distribution (blood 31\%,
    respiratory 23\%, feces 20\%, etc.) is systematically correlated with K-locus, our
    proportions are biased in an unknown direction. We papered over this by calling C1--C4
    ``partial (database growth)'' but the real issue is filter mismatch as much as growth.
  \item \textbf{Phylogenetic claim is proxied, not replicated.} Claim C20 (``9 clades'')
    was tested by counting 27 ST+K-locus combinations, which is not a phylogenetic tree.
    A real replication would run Roary + IQ-TREE/RAxML-ng on core-genome SNPs. This is
    doable on the same uicgpu box in $\sim$12\,h wall-clock, but was not run.
  \item \textbf{Differential-VF-gene count is not really tested.} C17 (35 differentially
    carried virulence genes, p$<$0.05) was verified only at the locus level (\emph{iuc},
    \emph{rmpA}, \emph{clb}) — not at the gene-by-gene level the paper reports. We report
    ``key loci confirmed different'' as PARTIAL, but the 35-gene count itself is
    untested; a fair reading is that C17 is closer to NOT TESTED than PARTIAL.
  \item \textbf{Tool-version mismatch is under-audited.} Paper uses Kleborate v2.x
    (2022); we use v3.2.4 (2026) with a rebuilt Kaptive database. We note this
    difference but did not quantify it (e.g., by running v2.x in parallel on 100 sampled
    isolates and computing the K-locus call concordance rate). Claim C13 (51 vs 19
    serotypes) may partly be a v2$\rightarrow$v3 taxonomic re-lumping artifact rather
    than a real reduction; we cannot tell without a controlled cross-version audit.
  \item \textbf{No assembly-quality filter.} All 955 assemblies were used regardless of
    N50 or completeness. Kleborate K-locus calls can misclassify on fragmented
    assemblies. See Open Question Q5.
  \item \textbf{Statistics missing.} The paper reports $\chi^2$ and Fisher's exact
    p-values throughout Tables 2 and 3. We reported percentages and percentage-point
    differences but did not compute the significance tests, so ``KL64 vscore = 2.06 vs.\
    KL47 = 1.80'' is not accompanied by a p-value. Directionally consistent, statistically
    unverified.
  \item \textbf{Publication provenance is not audited.} Both papers rely on BV-BRC public
    assemblies, whose depositors span hundreds of studies. Neither the paper nor this
    replication has audited whether a few large Chinese outbreak studies are inflating
    the KL64 signal via convenience sampling.
\end{itemize}

Read together: the epidemiological signal is real and reproducible; the mechanistic and
phylogenetic claims are asserted-consistent-with rather than tested; the ``REPLICATED''
verdict is appropriate for the descriptive-epidemiology backbone (C5--C11, C14, C16) but
should not be extended to the mechanism (C18) or the phylogenetic architecture (C20).

\section{Verdict}
\textbf{REPLICATED} for descriptive epidemiology, carbapenemase-gene distribution, and the
paper's central novel finding (KL47$\rightarrow$KL64 transition with virulence-gene
enrichment). \textbf{Not verified} for the proposed \emph{wzc}-recombination mechanism
(C18) and the phylogenetic-clade architecture (C20); those remain the highest-value
follow-on work. 18/20 claims tested (above the 80\% threshold); 0 contradicted; 17/18
tested claims (94\%) support the paper.

\section{Open Questions}
\begin{description}[leftmargin=*, style=nextline]
\item[Q1. Mechanism vs.\ clonal export.] What is the actual causal mechanism driving the
KL47$\rightarrow$KL64 transition --- \emph{wzc} CD1-VR2-CD2 recombination as the paper
proposes, plasmid-borne virulence selection, healthcare-network founder effects (one
super-spreader KL64 clone), or a combination?
\emph{Basis:} the paper explicitly names ``no experiment demonstrated the effect of
recombination'' as a limitation; our data cannot distinguish repeated conversion from
clonal export.
\emph{Next steps:} core-genome phylogeny + ancestral-state reconstruction on the 413
KL64 assemblies to count independent conversion events; lambda-Red \emph{wzc} allele
swap in a lab strain to directly test the recombination claim; search for KL64-\emph{wzc}
$+$ KL47-plasmid discordant isolates.

\item[Q2. Post-2020 trajectory.] Has KL64 dominance saturated, reversed, or been
overtaken (e.g., by KL24, KL105) since 2020, and are new carbapenemases (NDM-5,
OXA-232) replacing KPC-2 in these emerging serotypes?
\emph{Basis:} paper ends in 2020; our 2020 slice has n=14; KL24 (14.2\%) and KL105
(6.6\%) are non-negligible in our 955-genome pull and not emphasized by Zhang \emph{et al.}
\emph{Next steps:} extend the BV-BRC pull to 2021--2026, re-run Kleborate; cross-check
against WHO GLASS and CHINET; profile KL24/KL105 virulence-plasmid content.

\item[Q3. Metadata-selection bias.] How large is the selection bias between the paper's
386-genome curated set and a modern BV-BRC 955-genome pull, and does the
KL47$\rightarrow$KL64 finding survive stratification on country and sample source?
\emph{Basis:} 92\% missing \texttt{body\_sample\_site}; 61\% Chinese isolates on our
side vs 64.5\% on theirs; possible Chinese-hospital-network confounding.
\emph{Next steps:} recover sample-source metadata from NCBI BioSample via Entrez;
country/source-stratified Cochran-Mantel-Haenszel test; request Zhang \emph{et al.}'s
curated genome-inclusion table.

\item[Q4. Mobile-element vehicle for virulence.] Which specific MGEs (IS26, ICEKp10,
IncFIB/IncHI1B plasmids) are actually transporting \emph{rmpA/iuc} into ST11-KL64, and
is there a common vehicle that could be flagged for hv-CRKP surveillance?
\emph{Basis:} paper reports the differential presence but does not identify the vehicle;
Kleborate loci-level calls also do not resolve context; hv-CRKP has $>$50\% mortality
in published outbreaks.
\emph{Next steps:} long-read (Nanopore/PacBio) resequence a KL64-\emph{rmpA}$^{+}$
panel; run PlasmidFinder + mob\_suite + ISfinder around \emph{rmpA}; publish an
rmpA-vehicle typing scheme.

\item[Q5. Assembly-quality \& reassembly audit.] How much of the reported KL47/KL64
difference is real biology vs.\ artefacts of database growth, sequencer heterogeneity,
and assembly quality --- and does the finding reproduce blind on independently
re-assembled reads (SRA $\rightarrow$ de novo)?
\emph{Basis:} both papers use BV-BRC's pre-deposited assemblies with wildly variable
quality; Kleborate K-locus calls are sensitive to fragmentation; no assembly-QC filter
was applied.
\emph{Next steps:} compute N50/BUSCO for all 955 assemblies, filter, and re-run;
re-assemble 100 random SRA-available isolates from scratch (Unicycler on uicgpu); a
``no growth'' counterfactual restricted to the 386 pre-2022 accessions.
\end{description}

Full JSON versions of Q1--Q5 with detailed \texttt{next\_steps} are in
\texttt{report/open\_questions.json}.

\section{Reproducibility Envelope}
\begin{itemize}[leftmargin=*, itemsep=1pt]
  \item PDF source: MDPI open-access, CC-BY. Local copy \texttt{paper.pdf}
    (SHA-256 \texttt{a5b493...53df}, 1.52\,MB, 10 pages).
  \item Data source: BV-BRC REST API (public, no auth). Queries captured in \texttt{data/*.json}.
  \item Compute host: \texttt{uicgpu} (8$\times$A100 80\,GB, 2\,TB RAM), conda env
    \texttt{/data/stevens/envs/kleborate} (Python 3.10, Kleborate v3.2.4, Kaptive,
    AMRFinderPlus, Mash 2.3, minimap2 2.30).
  \item Wall-clock: BV-BRC pull $\sim$25\,min; assembly download $\sim$7\,min;
    Kleborate 8-way parallel $\sim$30\,min. Total $\sim$1\,h + human review.
  \item Cost: \$0 (uicgpu is free compute).
\end{itemize}

\end{document}
